% Included by main.tex; the preamble and references are shared.
\section[Reciprocal Green functions and a squared-distance kernel]
{Reciprocal Green functions\\and a squared-distance kernel}\label{sec:reciprocal-green}

Write $\mathbb D=\{z\in\mathbb C:|z|<1\}$ and $S^1=\partial\mathbb D$,
and identify $\mathbb C$ with $\mathbb R^2$. We use the solution kernel
$K$ and the boundary kernel $N$ constructed in \cref{sec:boundary-kernels}, under the
measure assumptions restated below.

A real symmetric kernel $H$ with zero diagonal is
\emph{conditionally negative semidefinite} if
\[
 \sum_{i,j=1}^n b_i b_j H(s_i,s_j)\le0
 \qquad\text{whenever }\sum_{i=1}^n b_i=0,
\]
for every finite family of points $s_i$ and real coefficients $b_i$.

\begin{samepage}
\begin{theorem}\label[theorem]{thm:squared-distance}
Let $dq=a\,dx$ on $\mathbb D$, where $a\in L^\infty(\mathbb D)$ is
nonnegative, and suppose that there is a constant $0<d<1$ such that
\[
 \int_{\mathbb D}|\psi|^2\,dq
 \le d\int_{\mathbb D}|\nabla\psi|^2\,dx
 \qquad(\psi\in H_0^1(\mathbb D)).
\]
For the associated kernels $K,N$ and every $p\in\mathbb D$, the kernel
\begin{equation}\label{eq:squared-distance}
 D_p(s,t)=\frac{K(p,s)K(p,t)}{N(s,t)}\quad(s\ne t),
 \qquad D_p(s,s)=0,
\end{equation}
is conditionally negative semidefinite on $S^1$.
\end{theorem}
\end{samepage}

\begin{remark}\label[remark]{rem:zero-potential}
When $q=0$, one has $K(p,s)=P_s(p)$ and
$N(s,t)=1/(\pi|s-t|^2)$, where
\[
 P_s(p)=\frac{1-|p|^2}{2\pi|s-p|^2}.
\]
The disk automorphism
$\phi_p(z)=(z-p)/(1-\overline p z)$ maps $p$ to zero and maps
$S^1$ onto itself. Direct subtraction gives
\[
 |\phi_p(s)-\phi_p(t)|^2
 =\frac{(1-|p|^2)^2|s-t|^2}{|s-p|^2|t-p|^2}
 \qquad(s,t\in S^1).
\]
Consequently
\[
 D_p(s,t)=\pi P_s(p)P_t(p)|s-t|^2
        =\frac1{4\pi}|\phi_p(s)-\phi_p(t)|^2.
\]
Thus $D_p$ is the squared Euclidean distance between the transformed
boundary points, multiplied by $1/(4\pi)$.
\end{remark}

For a real symmetric matrix $M$, let $n_+(M)$ denote its
\emph{positive index of inertia}: the number of strictly positive
eigenvalues, counted with multiplicity. For a real symmetric kernel $H$
on a set $X$, the assertion $n_+(H)\le1$ means that
\[
 n_+\bigl([H(x_i,x_j)]_{i,j=1}^n\bigr)\le1
\]
for every finite set of distinct points $x_1,\ldots,x_n\in X$.
It is a condition on finite sampled matrices. Likewise, a complex
kernel is positive semidefinite if every finite sampled matrix is
Hermitian positive semidefinite. Reciprocals and powers of kernels
below are entrywise, rather than operator inverses or compositions.

\begin{definition}[Reciprocal matrix condition]
\label[definition]{def:reciprocal-matrix-condition}
Let $H$ be a real symmetric kernel on $X$, with finite strictly
positive values off the diagonal and strictly positive, possibly
infinite, values on the diagonal. We say that $H$ satisfies the
\emph{reciprocal matrix condition} if, for every finite set of
distinct points $x_1,\ldots,x_n\in X$,
\[
 n_+(Q_H)\le1,\qquad
 (Q_H)_{ij}=\frac1{H(x_i,x_j)},
\]
where $1/(+\infty)=0$.
\end{definition}

The condition in \cref{def:reciprocal-matrix-condition} is related to the
McCullough--Quiggin characterization of complete Nevanlinna--Pick
kernels \cite{McC94,Qui94}. For an irreducible positive semidefinite kernel with finite,
nonzero entries, that characterization requires each finite reciprocal
matrix to have exactly one positive eigenvalue
\cite[Corollary 1.12]{AglerMcCarthy}. Here we use the weaker bound
$n_+\le1$, which also accommodates zero diagonal entries and limits.
For a reciprocal-matrix condition in an earlier Green-function setting,
see \cite[Section 2.2, Corollary 2.2.2 and Theorem 2.2.3]{Qui94}.
Invertible diagonal congruence preserves $n_+$ by Sylvester's law of
inertia. Subtracting a positive semidefinite matrix cannot increase
$n_+$: every subspace on which the difference is positive definite
is also a positive subspace for the original matrix. The bound
$n_+\le1$ also passes to principal submatrices and entrywise limits
of matrices of a fixed finite size.
For the latter assertion, use continuity of the ordered eigenvalues
of a real symmetric matrix with respect to its entries.

\begin{auxlemma}\label[auxlemma]{lem:rank-one-difference}
Let $B,C$ be real symmetric positive semidefinite $n\times n$
matrices, with $\operatorname{rank}C\le1$. Then
\[
 n_+(C-B)\le1.
\]
In particular, $n_+(uu^T-B)\le1$ for every $u\in\mathbb R^n$.
\end{auxlemma}
\begin{proof}
For $b\in\ker C$,
\[
 b^T(C-B)b=-b^TBb\le0.
\]
Since $\ker C$ has codimension at most one, every subspace of
dimension at least two meets it nontrivially. Thus $C-B$ cannot be
positive definite on a two-dimensional subspace, so it has at most
one positive eigenvalue. The particular case follows because
$uu^T$ is positive semidefinite and has rank at most one.
\end{proof}

\subsection{A bounded regularization of the Green kernel}

The Dirichlet Green function of the disk is
\[
 G(x,y)=\frac1{4\pi}\log\left(
  1+\frac{(1-|x|^2)(1-|y|^2)}{|x-y|^2}\right)
 \qquad(x\ne y),
\]
with $G(x,x)=+\infty$. A quadrature formula sampling $G$ at its
integration nodes would encounter these infinite diagonal values.
We therefore construct bounded replacements with finite positive
diagonal satisfying the condition in
\cref{def:reciprocal-matrix-condition}.

\begin{lemma}\label[lemma]{lem:green-regularization}
There are bounded continuous strictly positive kernels $C_\epsilon$
on $\mathbb D\times\mathbb D$, $0<\epsilon\le1$, such that
$C_\epsilon\uparrow G$ as $\epsilon\downarrow0$, including on the
diagonal, and each $C_\epsilon$ satisfies the reciprocal matrix
condition in \cref{def:reciprocal-matrix-condition}.
\end{lemma}
\begin{proof}
Define the two-variable kernel
\[
 h(x,y)=\frac{|x-y|^2}{(1-|x|^2)(1-|y|^2)}.
\]
For a scalar argument $\tau\ge0$, define
\[
 L(\tau)=\frac1{\log(1+1/\tau)}\quad(\tau>0),\qquad L(0)=0,
 \qquad f(\tau)=\tau+\frac12-L(\tau).
\]
Both scalar functions are continuous on $[0,\infty)$, and
\[
 \frac1{G(x,y)}=4\pi L(h(x,y)),
\]
with reciprocal zero on the diagonal.
We will show that $n_+(h+\tfrac12)\le1$ and that the kernel
$(x,y)\mapsto f(h(x,y))$ is positive semidefinite.

For a finite set $x_1,\ldots,x_n$, put
$A_{ij}=h(x_i,x_j)+\tfrac12$ and
$W=\operatorname{diag}(1-|x_i|^2)$. Expanding $|x_i-x_j|^2$ gives
\begin{equation}\label{eq:regularization-rank-one}
 (WAW)_{ij}
 =\frac12(1+|x_i|^2)(1+|x_j|^2)-2x_i\cdot x_j.
\end{equation}
The first term is the rank-one matrix $uu^T$, where
$u_i=(1+|x_i|^2)/\sqrt2$, and the matrix $[2x_i\cdot x_j]$
is a positive semidefinite Gram matrix. By
\cref{lem:rank-one-difference}, $n_+(WAW)\le1$.
Since $W$ is invertible, Sylvester's law of inertia gives $n_+(A)\le1$.

For $\tau>0$, direct integration gives
\begin{equation}\label{eq:logarithmic-mean-integral}
 \begin{aligned}
 \int_0^1\tau^{1-\alpha}(\tau+1)^\alpha\,d\alpha
 &=\tau\int_0^1
        \exp\bigl(\alpha\log(1+1/\tau)\bigr)\,d\alpha\\
 &=\frac{\tau\bigl((1+1/\tau)-1\bigr)}{\log(1+1/\tau)}
 =L(\tau).
 \end{aligned}
\end{equation}
For related integral representations of the weighted geometric mean
and the logarithmic mean, see \cite{QiZhangLi}. We derive the further
representation needed here.

For $0<\alpha<1$, let $\mathsf V_\alpha$ be a random variable with
probability density
\[
 c_\alpha\frac{v^{\alpha-1}}{1+v}\quad(v>0),
 \qquad
 c_\alpha=\left(\int_0^\infty
                      \frac{v^{\alpha-1}}{1+v}\,dv\right)^{-1}.
\]
The normalizing integral is finite at both zero and infinity. Set
$\mathsf R_\alpha=(1+\mathsf V_\alpha)^{-1}\in(0,1)$, and write
$\mathbb E_\alpha$ for expectation under this law. The sans-serif
symbols denote random variables; $v$ and $r$ denote fixed scalar
values. Integration by parts, with vanishing endpoint terms, gives
\[
 \mathbb E_\alpha(1-\mathsf R_\alpha)
 =c_\alpha\int_0^\infty\frac{v^\alpha}{(1+v)^2}\,dv
 =\alpha c_\alpha\int_0^\infty
                          \frac{v^{\alpha-1}}{1+v}\,dv
 =\alpha.
\]
For $\tau>0$, substituting $w=v\tau/(\tau+1)$ in the integral
below gives
\begin{equation}\label{eq:geometric-mean-integral}
 \begin{aligned}
 \tau^{1-\alpha}(\tau+1)^\alpha
 &=c_\alpha\int_0^\infty
   \frac{\tau(\tau+1)}{(\tau+1)+v\tau}\,v^{\alpha-1}\,dv\\
 &=\mathbb E_\alpha
       \frac{\tau(\tau+1)}{\tau+\mathsf R_\alpha}\\
 &=\mathbb E_\alpha\left[
       \tau+(1-\mathsf R_\alpha)
       -\frac{\mathsf R_\alpha(1-\mathsf R_\alpha)}
                    {\tau+\mathsf R_\alpha}\right].
 \end{aligned}
\end{equation}
At $\tau=0$ the last two expressions also agree, and both are zero.
Integrating in $\alpha$ and using
\cref{eq:logarithmic-mean-integral} therefore yields
\begin{equation}\label{eq:regularization-mixture}
 f(\tau)=\int_0^1\mathbb E_\alpha
       \frac{\mathsf R_\alpha(1-\mathsf R_\alpha)}
                    {\tau+\mathsf R_\alpha}\,d\alpha,
 \qquad 0\le f(\tau)\le f(0)=\frac12.
\end{equation}
The formula at zero follows directly from
$\mathbb E_\alpha(1-\mathsf R_\alpha)=\alpha$ and $f(0)=1/2$.
The bound follows from
$\mathsf R_\alpha/(\tau+\mathsf R_\alpha)\le1$.

We next prove the positive semidefiniteness needed in this mixture.
We use the standard Gram-kernel criterion and the Schur product
theorem \cite[Propositions 17.1 and 17.5]{GallierQuaintance}.
In particular, a pointwise convergent sum
$\sum_m\varphi_m(x)\overline{\varphi_m(y)}$ is positive
semidefinite: its finite partial sums are Gram kernels, and the
cone of positive semidefinite matrices is closed. More generally,
finite-valued pointwise limits of positive semidefinite kernels are
positive semidefinite; see also \cite[Problem 17.4]{GallierQuaintance}.

The power series
\[
 S(x,y)=\frac1{1-x\overline y}
       =\sum_{m=0}^\infty x^m\overline y^{\,m}
 \qquad(x,y\in\mathbb D)
\]
satisfies this criterion. Its complex conjugate is also positive
semidefinite, by the same criterion with $\varphi_m(x)=\overline x^{\,m}$.
The Schur product theorem shows that
$S(x,y)\overline{S(x,y)}=|1-x\overline y|^{-2}$ is positive
semidefinite. On every finite set, multiplying this kernel by
$(1-|x|^2)(1-|y|^2)$ is congruence by the positive diagonal matrix
$\operatorname{diag}(1-|x_i|^2)$ and preserves positive
semidefiniteness. Since
\[
 |1-x\overline y|^2
 =|x-y|^2+(1-|x|^2)(1-|y|^2),
\]
the resulting kernel is
\[
 k(x,y)=\frac1{1+h(x,y)}
       =\frac{(1-|x|^2)(1-|y|^2)}{|1-x\overline y|^2}.
\]
Thus $k$ is positive semidefinite and $0<k(x,y)\le1$.
For each fixed scalar $0<r<1$, the geometric series gives
\begin{equation}\label{eq:regularization-geometric-series}
 \frac1{h(x,y)+r}
 =\frac{k(x,y)}{1-(1-r)k(x,y)}
 =\sum_{m=0}^\infty(1-r)^m k(x,y)^{m+1}.
\end{equation}
Every positive integer entrywise power of $k$ is positive
semidefinite by repeated application of the Schur product theorem.
The coefficients in this series are nonnegative, so its partial
sums are positive semidefinite. The series converges everywhere,
including when $x=y$, because $(1-r)k(x,y)\le1-r<1$.
For each fixed finite set of points, the matrices of partial sums
converge entrywise. Their limit is positive semidefinite: for every
vector $b$, its quadratic form is the limit of the nonnegative
quadratic forms of the partial sums. Thus
$(x,y)\mapsto1/(h(x,y)+r)$ is positive semidefinite.

For any finite points $x_i$ and real coefficients $b_i$, the mixture
\cref{eq:regularization-mixture} now gives
\[
 \begin{aligned}
 \sum_{i,j}b_i b_j f(h(x_i,x_j))
 &=\int_0^1\mathbb E_\alpha\left[
   \mathsf R_\alpha(1-\mathsf R_\alpha)
   \sum_{i,j}\frac{b_i b_j}{h(x_i,x_j)+\mathsf R_\alpha}
   \right]\,d\alpha\\
 &\ge0.
 \end{aligned}
\]
For each realization of $\mathsf R_\alpha\in(0,1)$, the inner
quadratic form is nonnegative by the preceding argument.
The finite sum may be interchanged with the integrals because
\[
 0\le\frac{\mathsf R_\alpha(1-\mathsf R_\alpha)}
                {h(x_i,x_j)+\mathsf R_\alpha}
 \le1-\mathsf R_\alpha\le1.
\]
These integrands are measurable under the explicit densities above
and integrable in both $v$ and $\alpha$. This proves that
$(x,y)\mapsto f(h(x,y))$ is positive semidefinite, including on
the diagonal.

Define $C_\epsilon$ by
\begin{equation}\label{eq:regularized-green}
 \begin{aligned}
 \frac1{C_\epsilon(x,y)}
 &=4\pi\left[h(x,y)+\frac12
                  -(1-\epsilon)f(h(x,y))\right]\\
 &=4\pi\left[L(h(x,y))+\epsilon f(h(x,y))\right].
 \end{aligned}
\end{equation}
On any finite set, its reciprocal matrix is $4\pi$ times the
matrix of $h+\tfrac12$, minus $4\pi(1-\epsilon)$ times the matrix
of $f\circ h$. For the former matrix $A$,
\cref{eq:regularization-rank-one} shows that $WAW$ is a rank-one
positive semidefinite matrix minus a positive semidefinite Gram
matrix. Hence $n_+(WAW)\le1$ by \cref{lem:rank-one-difference}.
Since $W$ is invertible, Sylvester's
law of inertia gives $n_+(A)=n_+(WAW)\le1$.
The matrix of $f\circ h$ is positive semidefinite by the mixture
argument.
Since subtracting a positive semidefinite matrix cannot increase
$n_+$, the reciprocal matrix has at most one positive eigenvalue.

For $\tau\ge0$,
\[
 L(\tau)+\epsilon f(\tau)
 =\epsilon\left(\tau+\frac12\right)+(1-\epsilon)L(\tau)
 \ge\frac\epsilon2.
\]
The reciprocal in \cref{eq:regularized-green} is therefore at least
$2\pi\epsilon$. Continuity of $h,L,f$ shows that $C_\epsilon$
is continuous and strictly positive, with
$C_\epsilon\le1/(2\pi\epsilon)$.
Finally $f\ge0$ implies $C_\epsilon\uparrow G$ as
$\epsilon\downarrow0$ off the diagonal, and $0<C_\epsilon\le G$.
On the diagonal,
$C_\epsilon(x,x)=1/(2\pi\epsilon)\uparrow+\infty=G(x,x)$.
\end{proof}

\subsection{A finite resolvent update}

\begin{lemma}\label[lemma]{lem:finite-resolvent-update}
Let $M$ be a real symmetric $n\times n$ matrix with strictly
positive entries, and suppose that its entrywise reciprocal
$Q=[1/M_{ij}]_{i,j=1}^n$ satisfies $n_+(Q)\le1$.
For any index $z\in\{1,\ldots,n\}$ and any scalar $c\ge0$, define
\[
 \widetilde M_{ij}=M_{ij}+cM_{iz}M_{zj}.
\]
Then the entrywise reciprocal
$\widetilde Q=[1/\widetilde M_{ij}]_{i,j=1}^n$ also satisfies
$n_+(\widetilde Q)\le1$.
\end{lemma}
\begin{proof}
The case $c=0$ is immediate, since $\widetilde M=M$.
The assertion is also immediate when $n=1$. Hence assume
$c>0$ and $n\ge2$. Define the diagonal matrix
\[
 D=\operatorname{diag}(M_{1z},\ldots,M_{nz})
\]
and set $H=DQD$. All diagonal entries of $D$ are positive, so $D$
is invertible and Sylvester's law of inertia gives
$n_+(H)=n_+(Q)\le1$. By symmetry of $M$,
\[
 H_{ij}=\frac{M_{iz}M_{jz}}{M_{ij}}
       =\frac{M_{iz}M_{zj}}{M_{ij}}.
\]
Write $\gamma=M_{zz}>0$. The $z$-row and $z$-column of $H$ are
constant, because $H_{iz}=M_{iz}M_{zz}/M_{iz}=\gamma$.

After simultaneously reordering the rows and columns to place
index $z$ last, let
\[
 H_0=[H_{ij}]_{i,j\ne z}
\]
be the $(n-1)\times(n-1)$ principal submatrix obtained by deleting
the $z$-row and $z$-column. With $\mathbf1_{n-1}$ denoting the
all-ones vector of length $n-1$, the full matrix is
\[
 H=\begin{pmatrix}
 H_0&\gamma\mathbf1_{n-1}\\
 \gamma\mathbf1_{n-1}^T&\gamma
 \end{pmatrix}.
\]
Changing from the standard basis to
$e_1-e_n,\ldots,e_{n-1}-e_n,e_n$ gives an invertible congruence
transforming $H$ into
\[
 \begin{pmatrix}
 H_0-\gamma\mathbf1_{n-1}\mathbf1_{n-1}^T&0\\
 0&\gamma
 \end{pmatrix}.
\]
The last block already contributes one positive eigenvalue.
Since $n_+(H)\le1$, the first block must be negative semidefinite.
Consequently the full $n\times n$ matrix
\[
 U=\gamma\mathbf1_n\mathbf1_n^T-H
  =\begin{pmatrix}
    \gamma\mathbf1_{n-1}\mathbf1_{n-1}^T-H_0&0\\
    0&0
   \end{pmatrix}
\]
is positive semidefinite. Here $\mathbf1_n$ is the all-ones vector
of length $n$. Since $H_{ii}>0$,
\[
 0\le U_{ii}=\gamma-H_{ii}<\gamma.
\]
The Cauchy--Schwarz inequality for a positive semidefinite matrix
then gives
\[
 |U_{ij}|^2\le U_{ii}U_{jj}<\gamma^2,
 \qquad |U_{ij}|<\gamma.
\]

For the reciprocal matrix of the updated matrix $\widetilde M$,
use the same diagonal matrix $D$ and define
$\widetilde H=D\widetilde QD$. Its entries are
\[
 \widetilde H_{ij}
 =\frac{M_{iz}M_{zj}}{M_{ij}+cM_{iz}M_{zj}}
 =\frac{H_{ij}}{1+cH_{ij}}.
\]
Define the constant rank-one positive semidefinite matrix
\[
 C_c=\frac{\gamma}{1+c\gamma}\,
                    \mathbf1_n\mathbf1_n^T.
\]
Using $H_{ij}=\gamma-U_{ij}$, the entries of
$B_c=C_c-\widetilde H$ satisfy
\begin{equation}\label{eq:finite-update-series}
 \begin{aligned}
 (B_c)_{ij}
 &=\frac{\gamma}{1+c\gamma}
                -\frac{H_{ij}}{1+cH_{ij}}\\
 &=\frac{U_{ij}}
        {(1+c\gamma)(1+c\gamma-cU_{ij})}\\
 &=\sum_{m=0}^\infty
       \frac{c^m U_{ij}^{\,m+1}}{(1+c\gamma)^{m+2}}.
 \end{aligned}
\end{equation}
The series converges absolutely in every entry, because
$|cU_{ij}|<c\gamma<1+c\gamma$.
Each positive integer entrywise power of $U$ is positive semidefinite
by the Schur product theorem
\cite[Proposition 17.5]{GallierQuaintance}; this conclusion does not
require the entries of $U$ to be nonnegative. All coefficients
in \cref{eq:finite-update-series} are nonnegative, so the matrices
of partial sums are positive semidefinite. Their entrywise limit
$B_c$ is positive semidefinite, since each quadratic form of $B_c$
is the limit of the corresponding nonnegative quadratic forms.

Thus $\widetilde H=C_c-B_c$ is a rank-one positive semidefinite
matrix minus a positive semidefinite matrix.
By \cref{lem:rank-one-difference}, $n_+(\widetilde H)\le1$.
Since $D$ is invertible, Sylvester's law of inertia gives
$n_+(\widetilde Q)=n_+(\widetilde H)\le1$, as required.
\end{proof}

\subsection{The interior resolvent kernel}\label{sec:interior-resolvent}

Let $0\le\rho\le q$ be supported on a compact subset of $\mathbb D$.
Write $\langle f,g\rangle_\rho=\int fg\,d\rho$ and
$\|f\|_\rho=\langle f,f\rangle_\rho^{1/2}$ on the real space
$L^2(\rho)$. The Green operator
\[
 T_\rho f(x)=\int_{\mathbb D}G(x,y)f(y)\,d\rho(y)
\]
has $\|T_\rho\|\le d<1$, by the Green-operator estimate of \cref{sec:boundary-kernels}.
For each interior point $x$, the section $G(x,\cdot)$ belongs to
$L^2(\rho)$: the density of $\rho$ is bounded, and the logarithmic
singularity is locally square integrable. Define
\begin{equation}\label{eq:interior-resolvent-kernel}
 \begin{aligned}
 M_\rho(x,y)&=G(x,y)+
 \left\langle G(x,\cdot),(I-T_\rho)^{-1}G(y,\cdot)
 \right\rangle_\rho\quad(x\ne y),\\
 M_\rho(x,x)&=+\infty.
 \end{aligned}
\end{equation}

\begin{proposition}\label[proposition]{prop:interior-reciprocal}
The kernel $M_\rho$ satisfies the reciprocal matrix condition in
\cref{def:reciprocal-matrix-condition}. In particular, every finite
reciprocal matrix has zero diagonal and at most one positive eigenvalue.
\end{proposition}
\begin{proof}
First fix $0<\epsilon\le1$ and define the operator
\[
 T_{\rho,\epsilon}f(x)
 =\int_{\mathbb D}C_\epsilon(x,y)f(y)\,d\rho(y)
\]
on $L^2(\rho)$. The bounds $0<C_\epsilon\le G$ from
\cref{lem:green-regularization} give the pointwise inequality
\[
 |T_{\rho,\epsilon}f(x)|\le T_\rho|f|(x)
 \qquad\text{for $\rho$-almost every }x.
\]
Taking $L^2(\rho)$ norms gives
\[
 \|T_{\rho,\epsilon}f\|_\rho
 \le\|T_\rho|f|\|_\rho\le d\|f\|_\rho,
 \qquad \|T_{\rho,\epsilon}\|\le d<1.
\]
Hence $I-T_{\rho,\epsilon}$ is invertible. For every $x,y\in\mathbb D$,
including $x=y$, define
\begin{equation}\label{eq:regularized-interior-resolvent}
 \begin{aligned}
 M_{\rho,\epsilon}(x,y)&=C_\epsilon(x,y)\\
 &\quad+\left\langle C_\epsilon(x,\cdot),
 (I-T_{\rho,\epsilon})^{-1}C_\epsilon(y,\cdot)
 \right\rangle_\rho.
 \end{aligned}
\end{equation}
The sections in this formula belong to $L^2(\rho)$ because
$C_\epsilon$ is bounded and $\rho$ is finite. Thus both the kernel
sections and the operator in \cref{eq:interior-resolvent-kernel}
have been replaced by their $C_\epsilon$ counterparts.

We first prove the condition in \cref{def:reciprocal-matrix-condition} for
$M_{\rho,\epsilon}$ by quadrature. If $\rho=0$, then
$M_{\rho,\epsilon}=C_\epsilon$, and this follows directly from
\cref{lem:green-regularization}. For the quadrature argument assume
$\rho\ne0$, and fix a nonempty compact set $E\subset\mathbb D$
such that $\rho(\mathbb D\setminus E)=0$.
For each integer $m\ge1$, choose a finite Borel partition
\[
 E=\bigsqcup_{j=1}^{J_m}E_{m,j},
 \qquad \max_j\operatorname{diam}(E_{m,j})\le\frac1m,
\]
with nonempty cells. Such partitions are obtained by intersecting
$E$ with a grid of half-open squares of side $1/(\sqrt2\,m)$
and discarding empty intersections. For each
$j\in\{1,\ldots,J_m\}$ choose $z_{m,j}\in E_{m,j}$, and put
\[
 w_{m,j}=\rho(E_{m,j}),\qquad
 I_m=\{j:w_{m,j}>0\},\qquad
 \zeta_m(x)=z_{m,j}\quad(x\in E_{m,j}).
\]
The representatives are distinct because the cells are disjoint.
Only indices in $I_m$ contribute to quadrature sums, but choosing
representatives also in the zero-mass cells defines $\zeta_m$ on
all of $E$. In particular,
$\sup_{x\in E}|\zeta_m(x)-x|\le1/m$.

Define the quadrature kernel and its operator by
\[
 \begin{aligned}
 C_{\epsilon,m}(x,y)&=C_\epsilon(\zeta_m(x),\zeta_m(y))
                         &&(x,y\in E),\\
 T_{\rho,\epsilon,m}f(x)&=\int_E C_{\epsilon,m}(x,y)f(y)\,d\rho(y)
                         &&(x\in E).
 \end{aligned}
\]
Uniform continuity of $C_\epsilon$ on $E\times E$ gives
\begin{equation}\label{eq:quadrature-kernel-convergence}
 \delta_m=\sup_{x,y\in E}
 |C_{\epsilon,m}(x,y)-C_\epsilon(x,y)|\longrightarrow0.
\end{equation}
For $f\in L^2(\rho)$, Cauchy--Schwarz gives
\[
 |(T_{\rho,\epsilon,m}-T_{\rho,\epsilon})f(x)|
 \le\delta_m\int_E|f|\,d\rho
 \le\delta_m\sqrt{\rho(E)}\,\|f\|_\rho.
\]
Taking the $L^2(\rho)$ norm in $x$ yields the exact estimate
\begin{equation}\label{eq:quadrature-operator-convergence}
 \|T_{\rho,\epsilon,m}-T_{\rho,\epsilon}\|
 \le\rho(E)\delta_m\longrightarrow0.
\end{equation}
All operator norms in this argument are on $L^2(\rho)$ unless
another space is specified. Choose $m_0$ so that
$\rho(E)\delta_m\le(1-d)/2$ for $m\ge m_0$. Then
\begin{equation}\label{eq:quadrature-inverse-bound}
 \|T_{\rho,\epsilon,m}\|\le\frac{1+d}{2}<1,
 \qquad
 \|(I-T_{\rho,\epsilon,m})^{-1}\|\le\frac2{1-d}.
\end{equation}
The resolvent identity
\[
 \begin{aligned}
 &(I-T_{\rho,\epsilon,m})^{-1}-(I-T_{\rho,\epsilon})^{-1}\\
 &\quad=(I-T_{\rho,\epsilon,m})^{-1}
 (T_{\rho,\epsilon,m}-T_{\rho,\epsilon})
 (I-T_{\rho,\epsilon})^{-1}
 \end{aligned}
\]
therefore gives
\begin{equation}\label{eq:quadrature-resolvent-convergence}
 \left\|(I-T_{\rho,\epsilon,m})^{-1}
             -(I-T_{\rho,\epsilon})^{-1}\right\|
 \le\frac{2\rho(E)\delta_m}{(1-d)^2}\longrightarrow0.
\end{equation}

For each fixed evaluation point $x\in\mathbb D$, define
\[
 u_{m,x}(z)=C_\epsilon(x,\zeta_m(z)),\qquad
 u_x(z)=C_\epsilon(x,z)\quad(z\in E).
\]
Uniform continuity on $\{x\}\times E$ gives
\begin{equation}\label{eq:quadrature-source-convergence}
 \|u_{m,x}-u_x\|_\rho
 \le\sqrt{\rho(E)}\sup_{z\in E}
 |C_\epsilon(x,\zeta_m(z))-C_\epsilon(x,z)|\longrightarrow0.
\end{equation}
For $m\ge m_0$, define
\begin{equation}\label{eq:quadrature-resolvent-kernel}
 \begin{aligned}
 M_{\rho,\epsilon,m}(x,y)&=C_\epsilon(x,y)\\
 &\quad+\left\langle u_{m,x},
 (I-T_{\rho,\epsilon,m})^{-1}u_{m,y}\right\rangle_\rho.
 \end{aligned}
\end{equation}
The bounds and convergences in
\cref{eq:quadrature-inverse-bound,eq:quadrature-resolvent-convergence,eq:quadrature-source-convergence}
imply, for every fixed $x,y\in\mathbb D$,
\begin{equation}\label{eq:quadrature-evaluation-convergence}
 M_{\rho,\epsilon,m}(x,y)\longrightarrow M_{\rho,\epsilon}(x,y).
\end{equation}
Indeed, split the difference of the pairings by changing first
$u_{m,x}$, then the inverse, then $u_{m,y}$. Cauchy--Schwarz bounds
the resulting terms by the corresponding source and operator errors
times bounded norms. Thus \eqref{eq:quadrature-evaluation-convergence}
is pointwise convergence in the evaluation points.

We now express \cref{eq:quadrature-resolvent-kernel} as a finite
matrix formula. For this fixed $m$, let $N_m=|I_m|$, order the
positive-mass cells by increasing index, and define
\[
 e_{m,j}=\frac{\mathbf1_{E_{m,j}}}{\sqrt{w_{m,j}}}\quad(j\in I_m),
 \qquad
 \mathcal S_m=\operatorname{span}\{\mathbf1_{E_{m,j}}:j\in I_m\}
 \subset L^2(\rho).
\]
The subspace $\mathcal S_m$ consists of functions constant
$\rho$-almost everywhere on each partition cell. The vectors
$e_{m,j}$ form an orthonormal basis of this subspace, because
the cells are disjoint and $\|\mathbf1_{E_{m,j}}\|_\rho^2=w_{m,j}$.
Define the isometry
\[
 J_m:\mathbb R^{N_m}\to L^2(\rho),\qquad
 J_m b=\sum_{j\in I_m}b_j e_{m,j}.
\]
Its range is $\mathcal S_m$, and
$(J_m^*f)_j=w_{m,j}^{-1/2}\int_{E_{m,j}}f\,d\rho$.
Define the symmetric matrix $A_m$ and coordinate vectors
$\xi_{m,x}$ by
\[
 (A_m)_{ij}=\sqrt{w_{m,i}}\,
       C_\epsilon(z_{m,i},z_{m,j})\sqrt{w_{m,j}},
 \qquad
 (\xi_{m,x})_j=\sqrt{w_{m,j}}\,C_\epsilon(x,z_{m,j}).
\]
Direct substitution gives
\begin{equation}\label{eq:quadrature-matrix-representation}
 T_{\rho,\epsilon,m}=J_m A_m J_m^*,\qquad
 u_{m,x}=J_m\xi_{m,x}.
\end{equation}
Thus the operator is represented by $A_m$ on $\mathcal S_m$
and vanishes on $\mathcal S_m^\perp$. In particular,
\[
 \|A_m\|_{\mathbb R^{N_m}\to\mathbb R^{N_m}}
 =\|T_{\rho,\epsilon,m}\|\le\frac{1+d}{2}<1.
\]
Since $(I-T_{\rho,\epsilon,m})J_m=J_m(I-A_m)$,
\cref{eq:quadrature-matrix-representation} yields the exact formula
\begin{equation}\label{eq:quadrature-finite-resolvent}
 M_{\rho,\epsilon,m}(x,y)
 =C_\epsilon(x,y)+\xi_{m,x}^T(I-A_m)^{-1}\xi_{m,y}.
\end{equation}

The Neumann series $(I-A_m)^{-1}=\sum_{k=0}^\infty A_m^k$
converges in matrix operator norm. Expanding the matrix products gives
\begin{equation}\label{eq:quadrature-product-expansion}
 \begin{aligned}
 &\xi_{m,x}^T(I-A_m)^{-1}\xi_{m,y}\\
 &=\sum_{k=0}^\infty\sum_{j_0,\ldots,j_k\in I_m}
   \left(\prod_{\ell=0}^k w_{m,j_\ell}\right)
   C_\epsilon(x,z_{m,j_0})\\
 &\qquad\qquad{}\times\left(\prod_{\ell=0}^{k-1}
     C_\epsilon(z_{m,j_\ell},z_{m,j_{\ell+1}})\right)
   C_\epsilon(z_{m,j_k},y).
 \end{aligned}
\end{equation}
For $k=0$ the middle product is one. The intermediate indices may
repeat, and each occurrence contributes its quadrature weight.
All summands are nonnegative. Finiteness follows from
\eqref{eq:quadrature-finite-resolvent}, or directly from
\[
 |\xi_{m,x}^T A_m^k\xi_{m,y}|
 \le\|\xi_{m,x}\|\,\|\xi_{m,y}\|\,\|A_m\|^k.
\]

For a subset $J\subseteq I_m$, let $\mathcal R_{m,J}(x,y)$ denote
the series on the right-hand side of
\eqref{eq:quadrature-product-expansion} with every intermediate
index $j_0,\ldots,j_k$ restricted to $J$; put
$\mathcal R_{m,\varnothing}=0$. Nonnegativity gives
\begin{equation}\label{eq:restricted-series-bound}
 0\le\mathcal R_{m,J}(x,y)
 \le\mathcal R_{m,I_m}(x,y)
 =\xi_{m,x}^T(I-A_m)^{-1}\xi_{m,y}<\infty.
\end{equation}
Define
\[
 M^{(J)}(x,y)=C_\epsilon(x,y)+\mathcal R_{m,J}(x,y).
\]
In particular, $M^{(\varnothing)}=C_\epsilon$ and
$M^{(I_m)}=M_{\rho,\epsilon,m}$.

Choose $j_*\in I_m\setminus J$, and write
$z=z_{m,j_*}$ and $w=w_{m,j_*}>0$.
We calculate $M^{(J\cup\{j_*\})}$ by grouping the new terms
according to the number $k\ge1$ of occurrences of $j_*$ among
their intermediate indices. Splitting the sequence at these
occurrences leaves $k+1$ segments whose intermediate indices
belong to $J$. Summing over each segment gives a factor $M^{(J)}$;
its direct $C_\epsilon$ term accounts for an empty intermediate
sequence. Each of the $k$ occurrences contributes a factor $w$.
Therefore, by regrouping nonnegative series,
\begin{equation}\label{eq:additional-node-expansion}
 \begin{aligned}
 M^{(J\cup\{j_*\})}(x,y)
 &=M^{(J)}(x,y)\\
 &\quad+\sum_{k=1}^\infty w^k M^{(J)}(x,z)
       \bigl(M^{(J)}(z,z)\bigr)^{k-1}M^{(J)}(z,y).
 \end{aligned}
\end{equation}
Every $M^{(J)}$ value is finite and strictly positive by
\eqref{eq:restricted-series-bound}. The series in
\eqref{eq:additional-node-expansion} is finite by the same bound,
so its geometric ratio satisfies $wM^{(J)}(z,z)<1$.
Summing it gives
\begin{equation}\label{eq:additional-node-update}
 \begin{aligned}
 M^{(J\cup\{j_*\})}(x,y)
 &=M^{(J)}(x,y)+c_JM^{(J)}(x,z)M^{(J)}(z,y),\\
 c_J&=\frac{w}{1-wM^{(J)}(z,z)}>0.
 \end{aligned}
\end{equation}

Fix a finite set of distinct evaluation points, and form its union
with the positive-mass quadrature nodes, listed without repetitions.
Sample every $M^{(J)}$ on this same finite set. Order the indices
in $I_m$ as $j_1,\ldots,j_{N_m}$ and put
$J_0=\varnothing$, $J_\ell=\{j_1,\ldots,j_\ell\}$.
For $J_0$, the sampled matrix is that of $C_\epsilon$, whose
reciprocal has at most one positive eigenvalue by
\cref{lem:green-regularization}. At each step,
\eqref{eq:additional-node-update} is exactly the update in
\cref{lem:finite-resolvent-update}, with $z$ an index in the
sampled matrix. Applying that lemma successively for
$\ell=1,\ldots,N_m$ proves the same bound for the reciprocal
matrix of $M^{(I_m)}=M_{\rho,\epsilon,m}$.
Taking the principal submatrix on the evaluation points proves
that $M_{\rho,\epsilon,m}$ satisfies the condition in
\cref{def:reciprocal-matrix-condition}.

By \eqref{eq:quadrature-evaluation-convergence}, for any fixed
finite set $x_1,\ldots,x_n$,
$M_{\rho,\epsilon,m}(x_i,x_j)\to M_{\rho,\epsilon}(x_i,x_j)$.
The approximating kernels are at least $C_\epsilon$, by
\eqref{eq:quadrature-product-expansion}, so the limiting entries
are strictly positive. Hence their reciprocal matrices converge
entrywise as well. The bound $n_+\le1$ passes to this limit of
matrices of fixed size. Thus $M_{\rho,\epsilon}$ satisfies
the condition in \cref{def:reciprocal-matrix-condition} for each $\epsilon$.

Finally let $\epsilon\downarrow0$. For $k\ge0$, define
\begin{equation}\label{eq:regularized-iterated-integral}
 \begin{aligned}
 \mathcal I_{\rho,\epsilon,k}(x,y)
 &=\int_{E^{k+1}}C_\epsilon(x,z_0)\\
 &\quad{}\times\left(\prod_{\ell=0}^{k-1}C_\epsilon(z_\ell,z_{\ell+1})\right)
   C_\epsilon(z_k,y)\,
   d\rho(z_0)\cdots d\rho(z_k).
 \end{aligned}
\end{equation}
The middle product is one when $k=0$. If $\rho=0$, take any
nonempty compact $E\subset\mathbb D$; all these integrals are zero.
The operator Neumann series and the nonnegative integral kernels give
\begin{equation}\label{eq:regularized-correction-series}
 M_{\rho,\epsilon}(x,y)-C_\epsilon(x,y)
 =\sum_{k=0}^\infty\mathcal I_{\rho,\epsilon,k}(x,y).
\end{equation}
By \cref{lem:green-regularization}, $C_\epsilon\uparrow G$.
Thus each integrand in \eqref{eq:regularized-iterated-integral}
increases pointwise to the product with $G$ in place of every
$C_\epsilon$ factor. Monotone convergence for the integrals and
then their series gives
\begin{equation}\label{eq:green-correction-limit}
 \begin{aligned}
 \lim_{\epsilon\downarrow0}
 \sum_{k=0}^\infty\mathcal I_{\rho,\epsilon,k}(x,y)
 &=\sum_{k=0}^\infty\mathcal I_{\rho,G,k}(x,y)\\
 &=\left\langle G(x,\cdot),(I-T_\rho)^{-1}G(y,\cdot)
   \right\rangle_\rho,
 \end{aligned}
\end{equation}
where $\mathcal I_{\rho,G,k}$ is the same integral with $G$.
The last equality follows from the Neumann series for
$(I-T_\rho)^{-1}$; its $k$th pairing is precisely
$\mathcal I_{\rho,G,k}$.

The correction in \eqref{eq:green-correction-limit} is finite
for every $x,y\in\mathbb D$, including $x=y$, because
\begin{equation}\label{eq:green-correction-bound}
 \left|\left\langle G(x,\cdot),(I-T_\rho)^{-1}G(y,\cdot)
 \right\rangle_\rho\right|
 \le\frac{\|G(x,\cdot)\|_\rho\,\|G(y,\cdot)\|_\rho}{1-d}.
\end{equation}
For $x\ne y$, the direct term also has a finite limit, so
\[
 M_{\rho,\epsilon}(x,y)\longrightarrow M_\rho(x,y)<\infty.
\]
On the diagonal, the nonnegative correction gives
\[
 M_{\rho,\epsilon}(x,x)\ge C_\epsilon(x,x)
 =\frac1{2\pi\epsilon}\longrightarrow+\infty.
\]
Thus, for any fixed finite set of distinct points $x_1,\ldots,x_n$,
the reciprocal matrices converge entrywise to the matrix with
off-diagonal entries $1/M_\rho(x_i,x_j)$ and diagonal entries zero.
Each approximating matrix has at most one positive eigenvalue,
so the same is true of this limit. This is exactly
the condition in \cref{def:reciprocal-matrix-condition} for $M_\rho$.
\end{proof}

\subsection{The boundary limit}

We first record explicit integral formulas for the boundary kernels.
For any $0\le\rho\le q$ and $k\ge1$, define
\begin{equation}\label{eq:solution-iterated-integral}
 \begin{aligned}
 \mathcal K_{\rho,k}(p,s)
 &=\int_{\mathbb D^k}G(p,x_1)\\
 &\quad{}\times\left(\prod_{\ell=1}^{k-1}G(x_\ell,x_{\ell+1})\right)
 P_s(x_k)\,d\rho(x_1)\cdots d\rho(x_k).
 \end{aligned}
\end{equation}
For $k\ge0$, define
\begin{equation}\label{eq:interaction-iterated-integral}
 \begin{aligned}
 \mathcal R_{\rho,k}(s,t)
 &=\int_{\mathbb D^{k+1}}P_s(x_0)\\
 &\quad{}\times\left(\prod_{\ell=0}^{k-1}G(x_\ell,x_{\ell+1})\right)
 P_t(x_k)\,d\rho(x_0)\cdots d\rho(x_k).
 \end{aligned}
\end{equation}
Empty products in these formulas equal one. The nonnegative
integrals and series define
\begin{equation}\label{eq:potential-boundary-series}
 \begin{aligned}
 K_\rho(p,s)&=P_s(p)+\sum_{k=1}^\infty\mathcal K_{\rho,k}(p,s),\\
 R_\rho(s,t)&=\sum_{k=0}^\infty\mathcal R_{\rho,k}(s,t),\\
 N_\rho(s,t)&=N_0(s,t)+R_\rho(s,t),\qquad
 N_0(s,t)=\frac1{\pi|s-t|^2}\quad(s\ne t).
 \end{aligned}
\end{equation}
These are the definitions from \cref{sec:boundary-kernels} with $q$ replaced by $\rho$.
In particular $K_q=K$, $R_q=R$, and $N_q=N$. The integrals are
understood as nonnegative integrals before finiteness is asserted;
they do not require a Poisson section to belong to $L^2(q)$.

\begin{lemma}\label[lemma]{lem:boundary-reciprocal}
Let $s_1,\ldots,s_n\in S^1$ be distinct, and let $p\in\mathbb D$.
Define the real symmetric matrix
\begin{equation}\label{eq:boundary-bordered-matrix}
 \begin{gathered}
 \mathcal B=\begin{pmatrix}B&\ell\\\ell^T&0\end{pmatrix},\\
 B_{ij}=\begin{cases}
 1/N(s_i,s_j),&i\ne j,\\
 0,&i=j,
 \end{cases}
 \qquad \ell_i=\frac1{K(p,s_i)}.
 \end{gathered}
\end{equation}
Then $n_+(\mathcal B)\le1$.
\end{lemma}
\begin{proof}
First fix a compactly supported measure $0\le\rho\le q$.
Choose a nonempty compact set $E\subset\mathbb D$ with
$\rho(\mathbb D\setminus E)=0$. For $s\in S^1$ and $x\in\mathbb D$,
the Green formula gives, when $rs\ne x$,
\[
 G(rs,x)=\frac1{4\pi}\log\left(
 1+\frac{(1-r^2)(1-|x|^2)}{|rs-x|^2}\right).
\]
On any fixed compact interior set, $|rs-x|$ is bounded away from
zero for $r$ sufficiently close to one. Thus
$\omega_{r,s}(x)=(1-r^2)(1-|x|^2)/|rs-x|^2$ tends uniformly to
zero, and
\[
 \frac{G(rs,x)}{1-r}
 =\frac{(1+r)(1-|x|^2)}{4\pi|rs-x|^2}\cdot
   \frac{\log(1+\omega_{r,s}(x))}{\omega_{r,s}(x)}.
\]
Since $\log(1+u)/u\to1$ as $u\to0$, this proves
\begin{equation}\label{eq:green-poisson-limit}
 \frac{G(rs,x)}{1-r}\longrightarrow
 P_s(x)=\frac{1-|x|^2}{2\pi|s-x|^2},
\end{equation}
uniformly on every compact subset of $\mathbb D$.
For distinct $s,t\in S^1$, the same Green formula gives
\[
 G(rs,rt)=\frac1{4\pi}\log\left(
 1+\frac{(1-r^2)^2}{r^2|s-t|^2}\right).
\]
The fraction inside the logarithm tends to zero, and division by
$(1-r)^2$, using $(1-r^2)^2/(1-r)^2=(1+r)^2$, yields
\begin{equation}\label{eq:green-boundary-limit}
 \frac{G(rs,rt)}{(1-r)^2}\longrightarrow
 \frac1{\pi|s-t|^2}=N_0(s,t).
\end{equation}

Each Poisson section $P_s$ is bounded on $E$, so it belongs to
$L^2(\rho)$. The uniform convergence in
\eqref{eq:green-poisson-limit} gives
\begin{equation}\label{eq:poisson-section-convergence}
 \left\|\frac{G(rs,\cdot)}{1-r}-P_s\right\|_\rho
 \le\sqrt{\rho(E)}\sup_{x\in E}
 \left|\frac{G(rs,x)}{1-r}-P_s(x)\right|\longrightarrow0.
\end{equation}
Here $(I-T_\rho)^{-1}$ is a fixed bounded operator on $L^2(\rho)$,
with norm at most $1/(1-d)$. Expanding its Neumann series and
using \cref{eq:solution-iterated-integral,eq:interaction-iterated-integral,eq:potential-boundary-series}
gives the identities
\begin{equation}\label{eq:compact-boundary-pairings}
 \begin{aligned}
 R_\rho(s,t)&=\langle P_s,(I-T_\rho)^{-1}P_t\rangle_\rho,\\
 K_\rho(p,s)&=P_s(p)+
 \langle G(p,\cdot),(I-T_\rho)^{-1}P_s\rangle_\rho.
 \end{aligned}
\end{equation}
All sections in these pairings belong to $L^2(\rho)$, including
$G(p,\cdot)$ by the square integrability noted in \cref{sec:interior-resolvent}.

For $r>|p|$, the points $rs_1,\ldots,rs_n,p$ are distinct.
Writing out \eqref{eq:interior-resolvent-kernel} at these points gives
\begin{equation}\label{eq:scaled-resolvent-evaluations}
 \begin{aligned}
 \frac{M_\rho(rs_i,rs_j)}{(1-r)^2}
 &=\frac{G(rs_i,rs_j)}{(1-r)^2}\\
 &\quad+\left\langle\frac{G(rs_i,\cdot)}{1-r},
 (I-T_\rho)^{-1}\frac{G(rs_j,\cdot)}{1-r}\right\rangle_\rho
 \quad(i\ne j),\\
 \frac{M_\rho(p,rs_i)}{1-r}
 &=\frac{G(p,rs_i)}{1-r}\\
 &\quad+\left\langle G(p,\cdot),
 (I-T_\rho)^{-1}\frac{G(rs_i,\cdot)}{1-r}\right\rangle_\rho.
 \end{aligned}
\end{equation}
By \eqref{eq:poisson-section-convergence} and boundedness of the
resolvent, the first pairing tends to $R_\rho(s_i,s_j)$ and the
second tends to $K_\rho(p,s_i)-P_{s_i}(p)$, by
\eqref{eq:compact-boundary-pairings}. The direct terms tend to
$N_0(s_i,s_j)$ and $P_{s_i}(p)$ by
\cref{eq:green-poisson-limit,eq:green-boundary-limit}, using symmetry
of $G$ for the second term. Consequently
\begin{equation}\label{eq:boundary-resolvent-limits}
 \begin{aligned}
 \frac{M_\rho(rs_i,rs_j)}{(1-r)^2}&\longrightarrow N_\rho(s_i,s_j)
 \quad(i\ne j),\\
 \frac{M_\rho(p,rs_i)}{1-r}&\longrightarrow K_\rho(p,s_i).
 \end{aligned}
\end{equation}

Let $Q_{\rho,r}$ be the reciprocal matrix of $M_\rho$ sampled at
$rs_1,\ldots,rs_n,p$, with zero diagonal as in
\cref{def:reciprocal-matrix-condition}.
By \cref{prop:interior-reciprocal}, $n_+(Q_{\rho,r})\le1$.
Define
\[
 D_r=\operatorname{diag}(
 \underbrace{1-r,\ldots,1-r}_{n\text{ entries}},1).
\]
Since $D_r$ is invertible, Sylvester's law of inertia gives
$n_+(D_rQ_{\rho,r}D_r)\le1$. Its entries satisfy
\[
 \begin{aligned}
 (D_rQ_{\rho,r}D_r)_{ij}
 &=\frac{(1-r)^2}{M_\rho(rs_i,rs_j)}\quad(i\ne j,\ i,j\le n),\\
 (D_rQ_{\rho,r}D_r)_{i,n+1}
 &=\frac{1-r}{M_\rho(rs_i,p)}\quad(i\le n).
 \end{aligned}
\]
Every diagonal entry remains zero. The limits in
\eqref{eq:boundary-resolvent-limits} are strictly positive, since
$N_\rho\ge N_0>0$ off the diagonal and $K_\rho(p,s)\ge P_s(p)>0$.
Thus those limits and symmetry imply the entrywise matrix convergence
\begin{equation}\label{eq:compact-bordered-limit}
 \begin{gathered}
 D_rQ_{\rho,r}D_r\longrightarrow
 \mathcal B_\rho=\begin{pmatrix}B_\rho&\ell_\rho\\\ell_\rho^T&0\end{pmatrix},
 \\
 (B_\rho)_{ij}=\begin{cases}
 1/N_\rho(s_i,s_j),&i\ne j,\\0,&i=j,
 \end{cases}
 \quad (\ell_\rho)_i=\frac1{K_\rho(p,s_i)}.
 \end{gathered}
\end{equation}
Since the matrices have fixed size, the bound $n_+\le1$ passes to
the limit. Hence $n_+(\mathcal B_\rho)\le1$ for every compactly
supported $0\le\rho\le q$.

Now choose $0<r_1<r_2<\cdots\uparrow1$, set
$F_m=\{|x|\le r_m\}$, and define $d\rho_m=\mathbf1_{F_m}\,dq$.
To apply monotone convergence, write the integrals in
\cref{eq:solution-iterated-integral,eq:interaction-iterated-integral}
against the fixed product measures formed from $q$. Their
integrands are respectively the displayed products of $G$ and
$P_s$, multiplied by
\[
 \prod_{\ell=1}^k\mathbf1_{F_m}(x_\ell)
 \qquad\text{and}\qquad
 \prod_{\ell=0}^k\mathbf1_{F_m}(x_\ell).
\]
All factors are nonnegative. Since $F_m\uparrow\mathbb D$, these
indicator products increase pointwise to one. Monotone convergence
therefore gives, for each fixed $k$,
\[
 \mathcal K_{\rho_m,k}(p,s)\uparrow\mathcal K_{q,k}(p,s),
 \qquad
 \mathcal R_{\rho_m,k}(s,t)\uparrow\mathcal R_{q,k}(s,t).
\]
Applying monotone convergence to the nonnegative series in
\eqref{eq:potential-boundary-series} yields
\begin{equation}\label{eq:cutoff-boundary-kernel-limits}
 \begin{aligned}
 K_{\rho_m}(p,s)&\uparrow
 P_s(p)+\sum_{k=1}^\infty\mathcal K_{q,k}(p,s)=K(p,s),\\
 R_{\rho_m}(s,t)&\uparrow
 \sum_{k=0}^\infty\mathcal R_{q,k}(s,t)=R(s,t),\\
 N_{\rho_m}(s,t)&\uparrow N_0(s,t)+R(s,t)=N(s,t)\quad(s\ne t).
 \end{aligned}
\end{equation}
These limits are exactly the kernels constructed for $q$ in
\Cref{sec:boundary-kernels}. \Cref{lem:kernel-estimates} states that
$K(p,s)$ is finite for $p\in\mathbb D$ and $s\in S^1$, and that
$R(s,t)$ is finite for distinct $s,t\in S^1$. Thus the limiting
$K$ and $N$ values used here are finite and strictly positive.
Taking their reciprocals in \eqref{eq:cutoff-boundary-kernel-limits}
gives $\mathcal B_{\rho_m}\to\mathcal B$ entrywise, with the
diagonal identically zero. Since $n_+(\mathcal B_{\rho_m})\le1$,
the same bound holds for $\mathcal B$, proving the lemma.
\end{proof}

\begin{proof}[Proof of \cref{thm:squared-distance}]
For distinct $s_1,\ldots,s_n\in S^1$, use the matrix
$\mathcal B$ in \eqref{eq:boundary-bordered-matrix}.
By \cref{lem:boundary-reciprocal}, $n_+(\mathcal B)\le1$.
Each $K(p,s_i)$ is finite and strictly positive, so
$\ell_i=1/K(p,s_i)>0$ and $\ell\ne0$.

We first show that $v^TBv\le0$ whenever $\ell^Tv=0$.
Suppose instead that $\ell^Tv=0$ and $v^TBv>0$.
Set $y_0=\ell/(\ell^T\ell)$, so $\ell^Ty_0=1$, and define
\[
 \widehat y=y_0-\frac{v^TBy_0}{v^TBv}\,v.
\]
Then $v^TB\widehat y=0$ and $\ell^T\widehat y=1$.
Choose $c\in\mathbb R$ so that $\widehat y^TB\widehat y+2c>0$.
For the vectors $a_*=(v,0)$ and $b_*=(\widehat y,c)$,
\[
 \begin{aligned}
 a_*^T\mathcal B a_*&=v^TBv>0,\\
 b_*^T\mathcal B b_*&=\widehat y^TB\widehat y+2c>0,\\
 a_*^T\mathcal B b_*&=v^TB\widehat y+c\ell^Tv=0.
 \end{aligned}
\]
The vectors are linearly independent, because $\ell^Tv=0$ while
$\ell^T\widehat y=1$. The vanishing cross term shows that the
quadratic form of $\mathcal B$ is positive definite on their
two-dimensional span, contradicting $n_+(\mathcal B)\le1$.
Therefore
\begin{equation}\label{eq:bordered-form-negativity}
 v^TBv\le0\qquad\text{whenever }\ell^Tv=0.
\end{equation}

Now let $b_1,\ldots,b_n$ be real coefficients with $\sum_i b_i=0$,
and set $v_i=K(p,s_i)b_i$. Then
\[
 \ell^Tv=\sum_i\frac{K(p,s_i)b_i}{K(p,s_i)}=0,
\]
and, by \cref{eq:squared-distance,eq:boundary-bordered-matrix},
\[
 \sum_{i,j=1}^n b_i b_j D_p(s_i,s_j)=v^TBv\le0.
\]
The diagonal terms vanish on both sides, and the inequality follows
from \eqref{eq:bordered-form-negativity}. This proves conditional
negative semidefiniteness for distinct points. For a finite family
with repeated points, replace the coefficients at each distinct
point by their sum. This preserves both the total coefficient sum
and the quadratic expression, so the same conclusion holds.
\end{proof}
